\textbf{结论名称：}
双曲线焦点三角形面积公式

\textbf{适用条件：}
双曲线标准方程 \(\frac{x^2}{a^2}-\frac{y^2}{b^2}=1\)（或 \(\frac{y^2}{a^2}-\frac{x^2}{b^2}=1\)），\(P\) 为双曲线上一点，\(F_1,F_2\) 为两个焦点，\(\theta =\angle F_1PF_2\)（\(0<\theta<\pi\)），且 \(P\) 不在实轴端点。

\textbf{变量说明：}
\(a\) 为半实轴长，\(b\) 为半虚轴长，\(c=\sqrt{a^2+b^2}\) 为半焦距；\(S\) 为 \(\triangle F_1PF_2\) 的面积；\(\theta\) 为两焦半径 \(PF_1\) 与 \(PF_2\) 的夹角。

\textbf{核心公式：}
\[
\boxed{S_{\triangle F_1PF_2}=b^2\cot\frac{\theta}{2}}
\]

\textbf{使用场景：}
已知双曲线上一点与两焦点的夹角，直接求焦点三角形面积；或由面积反求夹角；在涉及焦点弦、离心率、渐近线的选填题中减少计算量。

\textbf{简单例子：}
双曲线 \(\frac{x^2}{4}-\frac{y^2}{3}=1\)，则 \(b^2=3\)。若 \(P\) 在右支且 \(\angle F_1PF_2=60^\circ\)，则
\[
S=3\cot\frac{60^\circ}{2}=3\cot30^\circ=3\sqrt{3}.
\]
（椭圆对应公式为 \(S=b^2\tan\frac{\theta}{2}\)，二者互为余切形式。）